% Multi-process and multithreaded web servers.  Each row of seven cells is
% the sequence of processing steps of one request.
% Usage: % Multi-process and multithreaded web servers.  Each row of seven cells is
% the sequence of processing steps of one request.
% Usage: % Multi-process and multithreaded web servers.  Each row of seven cells is
% the sequence of processing steps of one request.
% Usage: % Multi-process and multithreaded web servers.  Each row of seven cells is
% the sequence of processing steps of one request.
% Usage: \input{figures/l06-mp-mt}   (width ~116 mm)
\begin{tikzpicture}[x=1mm, y=1mm, font=\footnotesize\sffamily,
  proc/.style={draw, rounded corners=2pt, fill=white},
  cell/.style={draw, fill=ln-box, minimum width=3.6mm, minimum height=4mm,
    inner sep=0pt},
  cache/.style={draw, fill=ln-accent!15, align=center, inner sep=0.5pt,
    font=\scriptsize\sffamily},
  lab/.style={font=\scriptsize\sffamily},
  ttl/.style={font=\footnotesize\sffamily\bfseries, anchor=west},
  >={Stealth[length=1.4mm]}]
  % (a) multi-process.  The Japanese label of the cache box is wider than the
  % English one, so the cells start further left and the box sits further in.
  \iflangja{\def\lnMPcell{12.9}\def\lnMPcache{47.9}}%
           {\def\lnMPcell{13.8}\def\lnMPcache{48}}%
  \node[ttl] at (0,31) {\iflangja{(a) マルチプロセス}{(a) multi-process}};
  \foreach \y/\n in {22/1,12/2,2/3} {
    \draw[proc] (0,\y-4) rectangle (54,\y+4);
    \node[lab] at (5.5,\y) {P\n};
    \foreach \i in {0,...,6} { \node[cell] at (\lnMPcell+4.4*\i,\y) {}; }
    \node[cache, minimum width=11mm, minimum height=6mm] at (\lnMPcache,\y)
      {\iflangja{キャッシュ}{cache}};
  }
  \node[lab, text=ln-muted, anchor=north] at (27,-2.8)
    {\iflangja{プロセスごとに別のアドレス空間}{a separate address space per process}};
  % (b) multithreaded.  Same adjustment for the wider Japanese shared-state box.
  \iflangja{\def\lnMTsh{107.6}\def\lnMTlab{65.5}\def\lnMTcell{70.1}\def\lnMTarr{98.3}}%
           {\def\lnMTsh{109.5}\def\lnMTlab{66.5}\def\lnMTcell{73.8}\def\lnMTarr{102}}%
  \node[ttl] at (62,31) {\iflangja{(b) マルチスレッド}{(b) multithreaded}};
  \draw[proc] (62,-2) rectangle (116,26);
  \node[cache, minimum width=11mm, minimum height=17mm] (sh) at (\lnMTsh,12)
    {\iflangja{共有\\状態\\（キャッシュ）}{shared\\state\\(cache)}};
  \foreach \y/\n in {19/1,12/2,5/3} {
    \node[lab] at (\lnMTlab,\y) {T\n};
    \foreach \i in {0,...,6} { \node[cell] at (\lnMTcell+4.4*\i,\y) {}; }
    \draw[<->] (\lnMTarr,\y) -- (sh.west |- 0,\y);
  }
  \node[lab, text=ln-muted, anchor=north] at (89,-2.8)
    {\iflangja{1つのアドレス空間を共有}{one shared address space}};
\end{tikzpicture}
   (width ~116 mm)
\begin{tikzpicture}[x=1mm, y=1mm, font=\footnotesize\sffamily,
  proc/.style={draw, rounded corners=2pt, fill=white},
  cell/.style={draw, fill=ln-box, minimum width=3.6mm, minimum height=4mm,
    inner sep=0pt},
  cache/.style={draw, fill=ln-accent!15, align=center, inner sep=0.5pt,
    font=\scriptsize\sffamily},
  lab/.style={font=\scriptsize\sffamily},
  ttl/.style={font=\footnotesize\sffamily\bfseries, anchor=west},
  >={Stealth[length=1.4mm]}]
  % (a) multi-process.  The Japanese label of the cache box is wider than the
  % English one, so the cells start further left and the box sits further in.
  \iflangja{\def\lnMPcell{12.9}\def\lnMPcache{47.9}}%
           {\def\lnMPcell{13.8}\def\lnMPcache{48}}%
  \node[ttl] at (0,31) {\iflangja{(a) マルチプロセス}{(a) multi-process}};
  \foreach \y/\n in {22/1,12/2,2/3} {
    \draw[proc] (0,\y-4) rectangle (54,\y+4);
    \node[lab] at (5.5,\y) {P\n};
    \foreach \i in {0,...,6} { \node[cell] at (\lnMPcell+4.4*\i,\y) {}; }
    \node[cache, minimum width=11mm, minimum height=6mm] at (\lnMPcache,\y)
      {\iflangja{キャッシュ}{cache}};
  }
  \node[lab, text=ln-muted, anchor=north] at (27,-2.8)
    {\iflangja{プロセスごとに別のアドレス空間}{a separate address space per process}};
  % (b) multithreaded.  Same adjustment for the wider Japanese shared-state box.
  \iflangja{\def\lnMTsh{107.6}\def\lnMTlab{65.5}\def\lnMTcell{70.1}\def\lnMTarr{98.3}}%
           {\def\lnMTsh{109.5}\def\lnMTlab{66.5}\def\lnMTcell{73.8}\def\lnMTarr{102}}%
  \node[ttl] at (62,31) {\iflangja{(b) マルチスレッド}{(b) multithreaded}};
  \draw[proc] (62,-2) rectangle (116,26);
  \node[cache, minimum width=11mm, minimum height=17mm] (sh) at (\lnMTsh,12)
    {\iflangja{共有\\状態\\（キャッシュ）}{shared\\state\\(cache)}};
  \foreach \y/\n in {19/1,12/2,5/3} {
    \node[lab] at (\lnMTlab,\y) {T\n};
    \foreach \i in {0,...,6} { \node[cell] at (\lnMTcell+4.4*\i,\y) {}; }
    \draw[<->] (\lnMTarr,\y) -- (sh.west |- 0,\y);
  }
  \node[lab, text=ln-muted, anchor=north] at (89,-2.8)
    {\iflangja{1つのアドレス空間を共有}{one shared address space}};
\end{tikzpicture}
   (width ~116 mm)
\begin{tikzpicture}[x=1mm, y=1mm, font=\footnotesize\sffamily,
  proc/.style={draw, rounded corners=2pt, fill=white},
  cell/.style={draw, fill=ln-box, minimum width=3.6mm, minimum height=4mm,
    inner sep=0pt},
  cache/.style={draw, fill=ln-accent!15, align=center, inner sep=0.5pt,
    font=\scriptsize\sffamily},
  lab/.style={font=\scriptsize\sffamily},
  ttl/.style={font=\footnotesize\sffamily\bfseries, anchor=west},
  >={Stealth[length=1.4mm]}]
  % (a) multi-process.  The Japanese label of the cache box is wider than the
  % English one, so the cells start further left and the box sits further in.
  \iflangja{\def\lnMPcell{12.9}\def\lnMPcache{47.9}}%
           {\def\lnMPcell{13.8}\def\lnMPcache{48}}%
  \node[ttl] at (0,31) {\iflangja{(a) マルチプロセス}{(a) multi-process}};
  \foreach \y/\n in {22/1,12/2,2/3} {
    \draw[proc] (0,\y-4) rectangle (54,\y+4);
    \node[lab] at (5.5,\y) {P\n};
    \foreach \i in {0,...,6} { \node[cell] at (\lnMPcell+4.4*\i,\y) {}; }
    \node[cache, minimum width=11mm, minimum height=6mm] at (\lnMPcache,\y)
      {\iflangja{キャッシュ}{cache}};
  }
  \node[lab, text=ln-muted, anchor=north] at (27,-2.8)
    {\iflangja{プロセスごとに別のアドレス空間}{a separate address space per process}};
  % (b) multithreaded.  Same adjustment for the wider Japanese shared-state box.
  \iflangja{\def\lnMTsh{107.6}\def\lnMTlab{65.5}\def\lnMTcell{70.1}\def\lnMTarr{98.3}}%
           {\def\lnMTsh{109.5}\def\lnMTlab{66.5}\def\lnMTcell{73.8}\def\lnMTarr{102}}%
  \node[ttl] at (62,31) {\iflangja{(b) マルチスレッド}{(b) multithreaded}};
  \draw[proc] (62,-2) rectangle (116,26);
  \node[cache, minimum width=11mm, minimum height=17mm] (sh) at (\lnMTsh,12)
    {\iflangja{共有\\状態\\（キャッシュ）}{shared\\state\\(cache)}};
  \foreach \y/\n in {19/1,12/2,5/3} {
    \node[lab] at (\lnMTlab,\y) {T\n};
    \foreach \i in {0,...,6} { \node[cell] at (\lnMTcell+4.4*\i,\y) {}; }
    \draw[<->] (\lnMTarr,\y) -- (sh.west |- 0,\y);
  }
  \node[lab, text=ln-muted, anchor=north] at (89,-2.8)
    {\iflangja{1つのアドレス空間を共有}{one shared address space}};
\end{tikzpicture}
   (width ~116 mm)
\begin{tikzpicture}[x=1mm, y=1mm, font=\footnotesize\sffamily,
  proc/.style={draw, rounded corners=2pt, fill=white},
  cell/.style={draw, fill=ln-box, minimum width=3.6mm, minimum height=4mm,
    inner sep=0pt},
  cache/.style={draw, fill=ln-accent!15, align=center, inner sep=0.5pt,
    font=\scriptsize\sffamily},
  lab/.style={font=\scriptsize\sffamily},
  ttl/.style={font=\footnotesize\sffamily\bfseries, anchor=west},
  >={Stealth[length=1.4mm]}]
  % (a) multi-process.  The Japanese label of the cache box is wider than the
  % English one, so the cells start further left and the box sits further in.
  \iflangja{\def\lnMPcell{12.9}\def\lnMPcache{47.9}}%
           {\def\lnMPcell{13.8}\def\lnMPcache{48}}%
  \node[ttl] at (0,31) {\iflangja{(a) マルチプロセス}{(a) multi-process}};
  \foreach \y/\n in {22/1,12/2,2/3} {
    \draw[proc] (0,\y-4) rectangle (54,\y+4);
    \node[lab] at (5.5,\y) {P\n};
    \foreach \i in {0,...,6} { \node[cell] at (\lnMPcell+4.4*\i,\y) {}; }
    \node[cache, minimum width=11mm, minimum height=6mm] at (\lnMPcache,\y)
      {\iflangja{キャッシュ}{cache}};
  }
  \node[lab, text=ln-muted, anchor=north] at (27,-2.8)
    {\iflangja{プロセスごとに別のアドレス空間}{a separate address space per process}};
  % (b) multithreaded.  Same adjustment for the wider Japanese shared-state box.
  \iflangja{\def\lnMTsh{107.6}\def\lnMTlab{65.5}\def\lnMTcell{70.1}\def\lnMTarr{98.3}}%
           {\def\lnMTsh{109.5}\def\lnMTlab{66.5}\def\lnMTcell{73.8}\def\lnMTarr{102}}%
  \node[ttl] at (62,31) {\iflangja{(b) マルチスレッド}{(b) multithreaded}};
  \draw[proc] (62,-2) rectangle (116,26);
  \node[cache, minimum width=11mm, minimum height=17mm] (sh) at (\lnMTsh,12)
    {\iflangja{共有\\状態\\（キャッシュ）}{shared\\state\\(cache)}};
  \foreach \y/\n in {19/1,12/2,5/3} {
    \node[lab] at (\lnMTlab,\y) {T\n};
    \foreach \i in {0,...,6} { \node[cell] at (\lnMTcell+4.4*\i,\y) {}; }
    \draw[<->] (\lnMTarr,\y) -- (sh.west |- 0,\y);
  }
  \node[lab, text=ln-muted, anchor=north] at (89,-2.8)
    {\iflangja{1つのアドレス空間を共有}{one shared address space}};
\end{tikzpicture}
